\section{结论}

本文提出 \ourframework，一种面向自动化 MAS 设计的训练方法，它将编排表述为函数调用强化学习问题；我们还提出 \ourbenchmark，用于从五个维度对 MAS 与 SAS 进行受控比较。分析表明，MAS 的收益受到任务结构、验证协议和模型能力共同影响。在这些洞见的指导下，\ourframework 在多种公共基准上取得了稳定提升，并位于性能-成本帕累托前沿。

% providing a principled foundation for designing, evaluating, and deploying multi-agent systems.

% In this work, we advanced the study of automatic MAS design by contributing both a principled training framework and a controlled understanding of when multi-agent coordination is effective. We introduced \ourframework, a training-time approach that formulates MAS orchestration as a function-calling reinforcement learning problem, enabling holistic system-level reasoning while abstracting complex, goal-oriented sub-agents as callable functions. We also proposed \ourbenchmark, a controlled benchmark that characterizes tasks along five axes—\depth, \horizon, \breadth, \myparallel, and \robustness—allowing systematic comparison between single- and multi-agent systems. Through extensive analysis and experiments, we showed that the benefits of MAS are not universal but depend critically on task structure, verification protocols, and the capabilities of both the orchestrator and sub-agents. Experiments on diverse public benchmarks further validate the effectiveness of \ourframework under these insights. Together, \ourframework, \ourbenchmark, and our in-depth analysis provide a foundation for more principled design, evaluation, and deployment of multi-agent systems.

%\update{
%We advance automatic MAS design in both methodology and understanding. We propose \ourframework, a training-time approach that formulates MAS orchestration as a function-calling reinforcement learning problem, avoiding sequential, code-level execution and enabling complex, goal-oriented sub-agents to be abstracted as callable functions. We also introduce \ourbenchmark, a controlled benchmark that characterizes tasks along five axes: \depth, \horizon, \breadth, \myparallel, and \robustness. Using both, we conduct a rigorous analysis of when and why MAS are beneficial, showing that performance gains critically depend on task structure, evaluation protocols, and the capabilities of both the orchestrator and sub-agents, rather than holding universally. 

% Together, \ourframework, \ourbenchmark, and our analysis provide a more principled foundation for training and understanding multi-agent systems.
%}

% Looking ahread, we envision \ourbenchmark and \masbenchmark to serve as a strong basica and ebalable better funderting and and trianingb of MAS in the pursui of myulti-agent intellegence

% We proposed \ourframework, a
% novel training-time framework that formulates MAS orchestration as a function-calling reinforcement learning problem with holistic orchestration, generating an entire MAS at once. In MAS-Orchestra, complex, goal-oriented sub-agents are abstracted
% as callable functions, enabling global reasoning
% over system structure while hiding internal execution details. To rigorously study when and why MAS are beneficial, we introduce \ourbenchmark, a controlled benchmark that characterizes tasks
% along five axes: \depth, \horizon, \breadth, \myparallel, and \robustness. Our analysis reveals that
% MAS gains depend critically on task structure,
% evaluation protocols, and the capabilities of both
% orchestrator and sub-agents, rather than hold-
% ing universally. Guided by these insights, MAS-
% Orchestra achieves consistent improvements on
% public benchmarks including mathematical rea-
% soning, multi-hop QA, and search-based QA. To-
% gether, MAS-Orchestra and MASBENCH enable
% better training and understanding of MAS in the
% pursuit of multi-agent intelligence.

% This work advances automatic multi-agent system design along two complementary directions: scalable training and principled analysis. MAS-Orchestra shows that holistic, training-time orchestration enables global coordination without relying on sequential, code-level execution, while MASBENCH provides a controlled setting to study when such coordination is actually useful. Our findings demonstrate that multi-agent gains are not universal, but depend on task structure, evaluation protocols, and agent capabilities. Together, this framework and benchmark support more principled development and evaluation of multi-agent systems, moving MAS design from heuristic construction toward systematic understanding.


% \zixuan{TODO: Make sure the terms and plural are consistent}
